% ======================================================================
% SECTION 12: LIMITATIONS AND FUTURE WORK
% Two thematic blocks: per-bill limitations and future-work tracks.
% ======================================================================

[PROCESSING INSTRUCTION (Limitations block - Per-Bill and Cross-Bill
Limitations, 4 to 6 paragraphs): Read in full
- patients/paper/Deep-Research-3/part3_metrics_guardrails.md (the
  Open Questions and Limitations section)
- new-trial/national-24-7-trial/paper/full-paper/sections/limitations_future.tex
  (Track A vs Track B framing)
- patient-journey/README.md (the single-patient orientation note)
- sponsor/final_paper/README.md (the synthetic agent decision logic
  note)

For each of the seven proposed bills, write one short paragraph that
lists the most important limit by name. Use measured language; do
NOT overclaim. State explicitly:

(12.1.a) HR 4501 (Patient Self-Selection): The bill assumes
machine-readable eligibility data of high quality across all CMS-
billable EHRs. The ASTP/ONC hospital APIs brief documents that
adoption is uneven; the 12-month deadline is aggressive. The Mazor
2025 null result on AI notifications alone reminds us that
self-selection rights without a complete pathway from discovery to
consent to enrollment may not improve participation. Cite
\cite{astp_onc_hospital_apis_2024} and
\cite{mazor2025ai_notifications_trial}.

(12.1.b) HR 4502 (Robot-and-Humanoid Choice): The bill assumes that
qualified-robot USL scoring is comparable across vendors. The
unification/usl/ paper documents that USL scoring requires multi-site
calibration; humanoid robots in particular have a thinner clinical
evidence base than surgical robots. Cite \cite{kawchak_2026_18778219}.

(12.1.c) HR 4503 (Procedural Modification): The bill assumes that
real-time consent updates can be logged within 5 minutes by every
qualified Physical AI executor. The MCP-PAI infrastructure is
documented; field deployment is not. The bill's "modifiable steps"
matrix presumes IRB review at scale that has not been tested. Cite
\cite{kawchak_2026_18973368}.

(12.1.d) HR 4504 (Doctor and Nurse Error Reduction): The bill is the
most contentious of the seven because it touches workforce questions
beyond the project brief. The Huang 2025 null result indicates that
ML does not always exceed Cox regression on real-world structured
survival data; the SHIELD-RT prospective RCT is convincing for one
domain (radiotherapy AE prediction) but does not generalize without
prospective validation in each domain. Cite the four-simulation paper
\cite{kawchak_2026_19994945} and acknowledge the
TRIPOD+AI / CREMLS reporting gap. The bill respects human doctor and
nurse roles in empathy, consent counseling, and ethical oversight.

(12.1.e) HR 4505 (Real-Time Patient-Sponsor Direct Communication):
The bill assumes Paradigm-Health-style aggregator infrastructure
generalizes beyond TRAVERSE and STREAM-SCLC. The FDA RTCT pilot has
not yet released aggregate metrics; the 24-month deadline is
aggressive. Cite \cite{fda_rtct_press_2026}.

(12.1.f) HR 4506 (American Leadership): The bill is the most
ambitious; it assumes seven-year sustained federal investment which
is subject to appropriations cycles. The American Leadership Index
relies on year-over-year metric stability across 50 states. Cite
\cite{rank4_fdora_fda_modernization_2022}.

(12.1.g) HR 4507 (Patient Health Data Self-Custody): The bill
assumes EHR vendors will support a self-custody export endpoint at
zero cost to the patient within 12 months. The Cures Act
information-blocking enforcement record is mixed. Cite
\cite{federal_register_info_blocking_disincentives_2024}.

After the per-bill paragraphs, add one cross-cutting paragraph
stating:
- The bills are draft language by an independent author and have
  not been reviewed by Congress, FDA, or any state legislature.
- The bills are not endorsed by Anthropic, OpenAI, Google, or any
  technology vendor.
- The patient-control framework remains an operational hypothesis;
  outcome studies are required to confirm that the bills, taken
  together, increase participation and decrease error rates without
  unintended consequences.
- The bills do not address pediatric or vulnerable-population
  protections beyond the existing 21 CFR 50 Subpart D framework
  documented at national-platform/21cfr50_adapt/02_irb_review_pediatric.tex
  and at \cite{kawchak_2026_19040707}.]

[PROCESSING INSTRUCTION (Future Work block - Three Future Tracks, 3
to 4 paragraphs):

Track A - Single big model performing all bills' obligations: a
future generation of Claude Code Opus 4.7 (or equivalent foundation-
model agent) holds all seven bill texts plus the AVAILABLE
DIRECTORIES content in a single 1M-token context and emits
implementation deliverables (machine-readable eligibility reports,
USL transparency reports, modification matrices, error-rate
publications, real-time aggregator endpoints, American Leadership
Index, self-custody endpoints) in a single inference loop. This
track is the natural extension of the four-simulation paper at
new-trial/national-24-7-trial/paper/full-paper/. Cite
\cite{kawchak_2026_19994945} and \cite{claude_opus_47}.

Track B - Distributed agents per bill, coordinated by a smaller
orchestrator: each of the seven bills has its own dedicated agent
(self-selection agent for HR 4501, robot-choice agent for HR 4502,
modification agent for HR 4503, error-rate agent for HR 4504,
sponsor-channel agent for HR 4505, leadership-index agent for HR
4506, self-custody agent for HR 4507), all coordinated by a smaller
orchestrator that runs on hardware as constrained as the Core
i5-6200U with 4 GB RAM documented at
sponsor/final_paper/168_hours/instructions/core_i5_6200u_4gb/README.md.
This track is the natural extension of Simulation 4 at
sponsor/final_paper/168_hours/. Cite \cite{kawchak_2026_19396256}.

Track C - Patient-side native deployment: a future generation
deploys a single integrated patient-side agent that operates HR 4501
through HR 4507 from the patient's own device, using
patient-controlled credentials and HIPAA-compliant authentication
to interact with provider, payer, and sponsor endpoints. This track
addresses the patient priority thesis directly and is the natural
extension of patients/patient_robot_instructions_fixed.tex pages
1 through 10 plus the digital-twin infrastructure at
patient-journey/stage_03_digital_twin.py. Cite
\cite{kawchak_2026_18810541} and \cite{kawchak_2026_19119939}.

Close with one paragraph naming three concrete future deliverables:
(i) a TRIPOD+AI / CREMLS-compliant retrospective validation of HR
4504's error-rate baseline on a real-patient cohort, drawing on the
SHIELD-RT design pattern;
(ii) a public benchmark of HR 4501's machine-readable eligibility
report performance across the TrialGPT, PRISM, and TrialMatchAI
matching systems on a shared dataset;
(iii) an FDA-aligned RTCT pilot submission that uses Simulation 4
(sponsor/final_paper/168_hours/) as the sponsor-side input and HR
4505's patient-sponsor direct channel as the patient-side input.
Cite \cite{fda_rtct_press_2026},
\cite{kawchak_2026_19994945}, \cite{kawchak_2026_19396256}, and
\cite{kawchak_2026_19119939}.]
